\subsection{Background}

A recent study found that IoT traffic increased significantly, with overall data usage per device tripling between 2016 and 2018 \cite{Shafiq2013, Andrade2017}. Notably, downlink traffic increased sixfold, demonstrating the accelerated adoption and reliance on IoT devices \cite{Shafiq2013, Andrade2017}. Despite this growth, only 2\% of IoT devices are LTE-enabled, in striking contrast to the 41\% LTE adoption of non-IoT devices in Europe during the same period \cite{GSM2018}. Furthermore, about 30\% of IoT devices are still confined to 2G connectivity, highlighting the technology gap and the difficulty network operators face while switching to modern networks \cite{GSM2018}.

The movement patterns of IoT devices are diverse, with a substantial share remaining immobile. Approximately 30\% of devices connect to only one cellular tower per month, indicating their deployment in fixed applications such as smart meters or environmental sensors \cite{Shafiq2013}.

Interestingly, IoT devices have distinct traffic characteristics. Approximately 92\% of devices produce more uplink traffic than downlink traffic, highlighting the prominence of applications such as monitoring and telemetry that require devices to send more data than they receive \cite{Shafiq2013}. In businesses like manufacturing, the median uplink-to-downlink traffic ratio is 4.66, compared to a general ratio of 1.41 for IoT devices across all industries.

As IoT expands into more dynamic and high-mobility scenarios, such as vehicle-to-vehicle (V2V) communication, the challenges and requirements for data exchange become increasingly complex. Unlike traditional IoT applications, V2V communication involves highly dynamic environments where vehicles share real-time information to enhance road safety, optimize traffic flow, and enable cooperative driving. These scenarios demand robust communication systems capable of handling rapid changes in channel conditions caused by high relative velocities and multipath propagation effects. The distinctive characteristics of V2V networks emphasize the importance of designing communication protocols that account for the unique challenges posed by such high-mobility use cases.

The Doppler effect, caused by relative movement between vehicles, leads to frequency shifts that can distort signals and reduce the accuracy of data transmission. This phenomenon is particularly pronounced in scenarios with high relative velocities, such as highways, where vehicles can travel at speeds exceeding 200 km/h and may require data rates from 2 to 54Mbps \cite{7763324}. The resulting Doppler spectrum governs the variation in time of the channel gain, and its effects are only valid for short periods due to rapidly changing vehicular conditions. To address these challenges, various V2V channel models have been developed, including stochastic and geometry-based stochastic models \cite{Karedal2009}, which simulate real-world conditions by accounting for multipath propagation and scattering from surrounding objects.

Path loss modeling plays a vital role in understanding V2V communication under different environmental conditions \cite{Karedal2011}. For example, urban areas with street canyon effects introduce additional complexity, as signal reflections from buildings create multiple propagation paths. Advanced models, such as the two-ray and six-ray street canyon models \cite{7763324}, capture these effects to improve the accuracy of path loss estimation. These models are essential for designing robust communication systems capable of overcoming the challenges of highly dynamic vehicular environments

